\section*{Chapitre \ref{ch:s1:short-interest-borrow-and-crowding} --- Encours vendeur, emprunt et encombrement}

\begin{solution}{exo:s1:short-interest-borrow-and-crowding:1}
$0.1 \times 80 = 8$ millions d'actions vendues à découvert contre 0.8 million échangées par jour~: 10 jours de couverture.
\end{solution}

\begin{solution}{exo:s1:short-interest-borrow-and-crowding:2}
A prête un titre à B, qui le vend à découvert à C~; C prête ce même titre à D, qui le vend à découvert à E. Un seul titre a ainsi été vendu à découvert deux fois, et l'encours vendeur
en compte deux. Répéter l'opération et l'encours vendeur dépasse le flottant, comme les services de la SEC l'ont expliqué dans le cas de GameStop.
\end{solution}

\begin{solution}{exo:s1:short-interest-borrow-and-crowding:3}
À un taux d'utilisation de 80~\%, les frais valent $0.25~\% + 39.75~\% \times 0.5^3 = 5.22~\%$ par an. Sur 5~\% du capital pendant neuf mois~: $0.05 \times 5.22~\% \times 0.75 = 0.20~\%$ du
capital.
\end{solution}

\begin{solution}{exo:s1:short-interest-borrow-and-crowding:4}
Les frais sont les plus élevés là où le taux d'utilisation est le plus fort, ce qui désigne un ensemble de titres plus restreint et plus coûteux que les titres les plus vendus à découvert en proportion des actions~: le portefeuille
fondé sur les frais, côté vendeur, rapporte 7.9~\% par an en moyenne, contre 5.4~\% pour le portefeuille fondé sur l'encours vendeur. Les titres du portefeuille fondé sur les frais sous-performent d'un montant inférieur à leurs frais, et son
signal est plus faible (IC de 0.006 contre 0.061), car les frais sont plats pour la plupart des valeurs puis saturent.
\end{solution}

\begin{solution}{exo:s1:short-interest-borrow-and-crowding:5}
Elle ne verse aucuns frais, ne fait face ni à des rappels ni à des squeezes, et suppose seulement que les titres les plus vendus à découvert font en moyenne moins bien que les autres~; la version long--short doit
aussi vendre à découvert les titres où la vente à découvert est la plus coûteuse et la plus dangereuse.
\end{solution}

\begin{solution}{exo:s1:short-interest-borrow-and-crowding:6}
Chaque position vaut $0.5/20 = 2.5~\%$ du capital~; une hausse d'un facteur vingt représente un gain de 19 fois pour le titre, si bien que le portefeuille perd $2.5~\% \times 19 = 47.5~\%$ du capital.
\end{solution}

\begin{solution}{exo:s1:short-interest-borrow-and-crowding:7}
Les pertes médianes sont de 2.8~\% (un quart couvert en cinq jours), 4.0~\% (la moitié en cinq jours) et 5.7~\% (la moitié en dix jours). Selon la loi d'impact en racine carrée, le mouvement quotidien
croît avec la racine carrée de la part couverte~; répartir les mêmes achats sur davantage de jours réduit le mouvement de chaque jour mais, dans le scénario de stress du chapitre, le
cours continue de monter plus longtemps, de sorte que le mouvement cumulé augmente. Les squeezes réels sont pires~: d'autres acheteurs se joignent au mouvement, et les vendeurs à découvert confrontés à des pertes se couvrent plus vite.
\end{solution}

\begin{solution}{exo:s1:short-interest-borrow-and-crowding:8}
L'encours vendeur n'est connu qu'à la publication de la FINRA, plusieurs jours après la date de règlement-livraison~; l'utiliser à la date de règlement-livraison relève d'un biais d'anticipation. Le backtest doit aussi
imputer les frais d'emprunt et les rappels sur ses positions vendeuses, et présenter des résultats pondérés à égalité comme pondérés par la valeur, car l'effet est faible dans les portefeuilles pondérés par la valeur.
\end{solution}

\begin{solution}{pb:s1:short-interest-borrow-and-crowding:1}
\begin{enumerate}
\item 122.97~\% du flottant en janvier 2021~; lorsque les actions vendues à découvert sont à nouveau prêtées par leurs acheteurs.
\item D'un cours de clôture de 19.95~\$ le 12 janvier à 347.51~\$ le 27 janvier (plus de 1\,600~\% au-dessus du 11 janvier), un plus haut intrajournalier de 483.00~\$ le 28 janvier, et
environ 2\,700~\% entre le plus bas du 8 janvier et ce plus haut.
\item Les rachats de couverture des principaux vendeurs à découvert ont coïncidé avec certaines fortes hausses, mais ils ne représentaient qu'une faible part des achats, et les cours sont restés élevés après leur reflux.
\item Deux fois par mois selon la Rule 4560 de la FINRA, à échéance de deux jours ouvrés après la date de règlement-livraison et publiés plus tard~: des données vieilles de plusieurs semaines au moment de leur utilisation.
\item Les titres fortement vendus à découvert ont sous-performé de 1.16~\% ajusté du risque sur 20 jours (2000--2004)~; beaucoup de régularités ne sont pas robustes, et les portefeuilles pondérés par la valeur à fort
encours vendeur n'ont pas sous-performé.
\item Offre prêtable d'environ 20~\% des actions~; des vendeurs informés suivent le rendement espéré injecté et des vendeurs non informés apportent du bruit~; des frais qui montent au-delà de 60~\%
d'utilisation, jusqu'à 40~\%~; rappels plus probables à fort taux d'utilisation.
\item 0.061, 0.048, 0.006 et 0.058.
\item 8.3~\% par an brut, 2.7~\% de frais et 0.95~\% de coûts~; ratios de Sharpe de 3.33, 2.24 et 1.84.
\item Ses titres sous-performent d'un montant inférieur à leurs frais (7.9~\% par an du côté vendeur).
\item 0.94~\% par an au-dessus du marché pondéré à égalité, $t = 7.2$, sans frais.
\item Les titres présentant un risque de vente à découvert plus élevé (frais susceptibles de bondir, prêts susceptibles d'être rappelés) ont des rendements plus faibles, une efficience des prix moindre et moins de
ventes à découvert.
\item Une probabilité quotidienne qui passe de 0.05~\% à 2~\% lorsque le taux d'utilisation dépasse 80~\%~: un fort taux d'utilisation est précisément là où se trouvent les titres fortement vendus à découvert.
\item \textbf{Résultat phare.} Le portefeuille par déciles d'encours vendeur rapporte 8.3~\% par an avant frais d'emprunt et 4.6~\% après frais et coûts (ratio de Sharpe de 3.33, 2.24
après frais)~; un portefeuille vendeur sur les vingt titres les plus encombrés, avec la moitié de son capital, perd une médiane de 4.0~\% du capital et, au pire, 8.3~\%, si la moitié de leur encours vendeur
se couvre en moins de cinq jours, contre une pire perte observée sur cinq jours de 3.0~\%.
\item La loi d'impact en racine carrée est ajustée sur des ordres ordinaires, et non sur des jours d'achat représentant plusieurs fois le volume, lorsque d'autres acheteurs se joignent au mouvement.
\item Elle coûte 47.5~\% du capital.
\item Plafonner chaque titre selon ses jours de couverture et l'exposition globale à l'encombrement du portefeuille, et le soumettre à des squeezes plutôt qu'à l'historique.
\item Évitement.
\item Les données quotidiennes de prêt de titres (taux d'utilisation, frais, rappels) et le volume quotidien des ventes à découvert, au-delà de l'encours vendeur publié deux fois par mois.
\item La queue d'un portefeuille vendeur se situe hors de la plupart des historiques~; elle doit être imposée par des scénarios de stress.
\item Les frais d'emprunt, le risque de rappel et une perte sans borne.
\end{enumerate}
\end{solution}

\begin{solution}{iq:s1:short-interest-borrow-and-crowding:1}
Les opérateurs informés anticipent une baisse, en moyenne, et la vente à découvert est coûteuse et risquée~; un encours vendeur élevé signifie aussi un achat potentiellement forcé
si les vendeurs doivent se couvrir.
\end{solution}

\begin{solution}{iq:s1:short-interest-borrow-and-crowding:2}
Des frais d'emprunt qui peuvent monter, des prêts qui peuvent être rappelés, des appels de marge, des pertes sans borne, et des squeezes lorsque de nombreux vendeurs doivent acheter en même temps.
\end{solution}

\begin{solution}{iq:s1:short-interest-borrow-and-crowding:3}
Noter chaque vente à découvert selon l'encours vendeur, le taux d'utilisation et les jours de couverture~; limiter les positions selon les jours de couverture et agréger l'exposition encombrée du portefeuille~; soumettre le portefeuille
à des scénarios de squeeze~; suivre les frais et les rappels au quotidien.
\end{solution}

\begin{solution}{iq:s1:short-interest-borrow-and-crowding:4}
Parce que la vente à découvert coûte des frais qui peuvent dépasser la sous-performance, et comporte un risque de rappel et un \omterm{def:s1:short-interest-borrow-and-crowding:squeeze}{risque de squeeze} qui absorbe la moyenne~; le signal est mieux utilisé
pour éviter ces titres.
\end{solution}

\begin{solution}{iq:s1:short-interest-borrow-and-crowding:5}
La date de règlement-livraison de chaque rapport sur l'encours vendeur et sa date de publication~; les données de frais et de taux d'utilisation datées au moment de leur observation~; les fenêtres de volume alignées sur elles.
Sans elles, le backtest utilise une information antérieure à sa publication.
\end{solution}

\begin{solution}{iq:s1:short-interest-borrow-and-crowding:6}
Le mouvement quotidien vaut environ $\eta\sigma\sqrt{Q/(dV)}$, de sorte que sur $d$ jours le cours monte d'environ $d\,\eta\sigma\sqrt{Q/(dV)} = \eta\sigma\sqrt{dQ/V}$ si
l'impact persiste. Il sous-estime un squeeze, car la loi est estimée sur des ordres bien plus petits que le volume quotidien, et parce qu'un squeeze s'auto-alimente~:
la hausse des cours force davantage de rachats de couverture et attire des acheteurs de momentum, ce que la loi ne contient pas.
\end{solution}
